\section{Embedded Code}\label{sec:EmbeddedCode}
Sometimes, it is necessary to embed code from another (programming) language into the \Java source code\footnote{If you write a code generator tool, this could even by Java code.}.

First, the code should be placed into a text block\index{text block} (see \autocite{ORACLE_DOC_LANGUAGE_SPECIFICATION:TextBlocks} for the definition):
\begin{lstlisting}[numbers=left]
final var code = """
    <code1>
        <code2>
    """;
\end{lstlisting}

The left margin of the text block is indicated by the position of the closing three double quotes. In the sample above, \verb#<code1># in line~2 is not indented at all, while \verb#<code2># in line~3 is indented by four spaces. The indentation\index{indentation} within the source code does not matter here; it is just important for optical reasons.

Aligning the opening and the closing delimiters (the three double quotes) is not recommended. This requires reindentation of the text block\index{text block} if the variable name or modifiers are changed.

Although I requested in \tqfullvref{sec:IndentationCharacter} to use only blanks for the indentation\index{indentation!character}, it is possible that the embedded code requires tabs:
\begin{lstlisting}[numbers=left]
final var codeJavaScript = """
    \t	setTimeout(function() {
    \t	\t	myDisplayer("Hello!");
    \t	}, 3000);
    """;
\end{lstlisting}
But for the indentation\index{indentation} of the text block\index{text block} itself (the indentation for the closing delimiter), you still use the blanks. Be careful to not mix the indentation characters\index{indentation!character} within the text block\index{text block} itself. It could lead to funny results.

\autocite{Laskey:TextBlocks} is a comprehensive guide to text blocks\index{text block} in \Java.

\subsubsection{Principles}
\begin{enumerate}[label=P\arabic*.]
\item{The opening three double-quotes (the \bsq{\textit{delimiter}}) for a text block\index{text block} should be placed at the right end of the previous line, and the closing delimiter should be placed on its own line, at the left margin of the text block.

Use the \lstinline|\| (line-terminator) escape sequence if the trailing line-feed is not wanted.}
\item{Avoid aligning the opening and closing delimiters and the text block's left margin.}
\item{Either use only spaces or only tabs for the indentation\index{indentation} inside a text block. Mixing white space will lead to a result with irregular indentation.}
\end{enumerate}

%------------------------------------------------------------------------------

\subsection{Formatting SQL inside Java}\label{sec:FormattingSQLInsideJava}

%------------------------------------------------------------------------------

\subsection{Formatting XML inside Java}\label{sec:FormattingXMLInsideJava}
You should embed only small fragments of an XML document into the Java source; larger fragments and full documents can be handled better when provided as resources.

%------------------------------------------------------------------------------

\subsection{Formatting HTML inside Java}\label{sec:FormattingHTMLInsideJava}
Same as for XML, also only small HTML fragments should be embedded into the Java source code. Anything else should go into a resource file.


%==============================================================================
% End of file!!
%==============================================================================
